\chapter*{Introducción}

\addcontentsline{toc}{chapter}{Introducción}

Los residuos de aparatos eléctricos y electrónicos (RAEE) combinan un alto impacto ambiental con un importante potencial económico, al contener materiales críticos y metales de valor que los convierten en una fuente secundaria estratégica de recursos \cite{leeEnvironmentalEconomicImpacts2024b}, lo que ha favorecido su integración en enfoques de economía circular \cite{heDrivingSustainableCircular2024}. Sin embargo, la transición hacia sistemas efectivos de aprovechamiento enfrenta una limitación recurrente: la dificultad de clasificar los residuos de manera eficiente en etapas tempranas, condición que determina la calidad de los materiales recuperados y la viabilidad de los procesos posteriores \cite{janoClosingLoopSystematic2026}.\\

La inteligencia artificial, en particular la visión computacional basada en aprendizaje profundo, ha mostrado una notable capacidad para reconocer patrones visuales complejos y discriminar entre tipos de residuos \cite{sharmaEnablersComputerVision2023}. Arquitecturas como las redes neuronales convolucionales y los modelos YOLO destacan por su desempeño en clasificación automatizada \cite{rajeevAdvancingEwasteClassification2025a}, y los enfoques orientados a la identificación por tipo de objeto (placas, baterías, cables o dispositivos) permiten estructurar salidas interpretables y aplicables a decisiones operativas \cite{ilievEmbeddedSystemEwaste2023}. La integración de estos modelos con sistemas físicos de manipulación ha mejorado la eficiencia y consistencia de los procesos de separación \cite{wangImprovedYOLObasedWaste2023}, mientras que arquitecturas más ligeras optimizan el desempeño sin aumentar los costos computacionales \cite{oiseEViTNetLightweightVision2026}.\\

No obstante, la mayoría de estos desarrollos se han concentrado fuera de América Latina. En Colombia, la gestión de RAEE se caracteriza por dinámicas informales, baja trazabilidad y escasa integración tecnológica en los procesos de clasificación; las iniciativas nacionales se han enfocado en aspectos normativos, logísticos y de sensibilización, sin abordar el desarrollo de soluciones tecnológicas aplicadas a la automatización de esta etapa \cite{guiaRAEEColombia2023}. La literatura revisada no evidencia implementaciones de modelos de visión computacional orientados a clasificar RAEE por tipo de objeto en el contexto nacional, lo que representa una oportunidad clara de investigación.\\

Dado que la clasificación constituye un punto de entrada para mejorar la eficiencia del sistema y generar información estructurada sobre los residuos (insumo fundamental para las estrategias de aprovechamiento dentro de la economía circular), esta investigación propone el desarrollo de un modelo de inteligencia artificial basado en visión computacional para la clasificación de residuos electrónicos, como aporte a la gestión de RAEE en el contexto colombiano.